\documentclass{sysuredhead}

\RedHeadSetup{
  layout = campus,
  direction = parallel,
  authority = {示例大学社会科学学院},
  mark = {《形势与政策》课程报告},
  document-number = {课程〔2026〕16号},
  title = {基层公共服务数字化的观察与思考},
  recipients = {课程教学组：},
  issuer = {示例大学社会科学学院学生},
  date = {2026年9月2日},
  note = {本文为虚构课程作业。},
  copies = {课程教学组。},
  print-office = {示例大学社会科学学院},
  print-date = {2026年9月2日},
  draft-label = {校园创作示例，不具公文效力},
  seal = none
}

\begin{document}

\RedHeadMakeHead

本文通过虚构观察材料，说明如何在保留校园创作表达的同时复用统一的公文排版组件。

\section{观察方法}

围绕信息获取、事项理解和意见反馈三个环节设计观察表，不记录真实个人信息。

\subsection{材料范围}

仅使用公开资料和自行编写的测试材料。

\section{主要认识}

数字化不只是将线下表格搬到线上，还需要重新梳理信息结构和服务流程。

\begin{RedHeadAttachments}
  \item 观察表样例
\end{RedHeadAttachments}

\RedHeadMakeClosing

\begin{RedHeadAttachment}{1}{观察表样例}
本附件为校园报告附件样式示例。
\end{RedHeadAttachment}

\RedHeadMakeImprint

\end{document}
